%!TEX program = xelatex
\documentclass[10.5pt, a4paper]{article}

\usepackage{fontspec}
\usepackage{xeCJK}
\usepackage{geometry}
\usepackage{setspace}
\usepackage{titlesec}
\usepackage{amsmath}
\usepackage{amssymb}
\usepackage{amsthm}
\usepackage{xcolor}
\usepackage{fancyhdr}
\usepackage{tcolorbox}

\geometry{top=1.6cm, bottom=1.6cm, left=2.0cm, right=2.0cm}
\setlength{\headheight}{14pt}

\setCJKmainfont{Microsoft JhengHei}
\setmainfont{Times New Roman}

\definecolor{accent}{RGB}{30, 80, 160}
\definecolor{thmfill}{RGB}{237, 246, 255}
\definecolor{warnfill}{RGB}{255, 243, 235}

\tcbuselibrary{breakable}

\pagestyle{fancy}
\fancyhf{}
\renewcommand{\headrulewidth}{0.4pt}
\fancyhead[R]{\small\color{gray} 黃崇晉 | L16141149 | 數論（一）2026}
\fancyfoot[C]{\small\thepage}

\titleformat{\section}[block]
  {\normalsize\bfseries\color{accent}}{}
  {0pt}{}
  [\vspace{1pt}\titlerule\vspace{3pt}]

\setstretch{1.22}
\setlength{\parskip}{2.5pt}
\setlength{\parindent}{0pt}

% compact list
\setlength{\topsep}{2pt}
\setlength{\partopsep}{0pt}
\makeatletter
\let\orig@itemize\itemize
\renewcommand{\itemize}{%
  \orig@itemize\setlength{\itemsep}{1pt}\setlength{\parsep}{0pt}%
}
\makeatother

% theorem environments
\newtheoremstyle{cpt}{2pt}{2pt}{\normalfont}{0pt}{\bfseries\color{accent}}{.}{0.45em}{}
\theoremstyle{cpt}
\newtheorem{definition}{定義}
\newtheorem{lemma}{引理}
\newtheorem{theorem}{定理}
\newtheorem{corollary}{推論}

% proof in a box
\newenvironment{proofbox}[1][證明]{%
  \begin{tcolorbox}[colback=thmfill, colframe=accent!40,
    breakable, left=5pt, right=5pt, top=3pt, bottom=3pt, boxrule=0.4pt]
  \noindent\textbf{\color{accent}#1.}\quad\small
}{%
  \hfill$\square$\end{tcolorbox}\vspace{-1pt}
}

\begin{document}

% ── Title ─────────────────────────────────────────────────────────────
\begin{center}
  {\LARGE\bfseries\color{accent}
    Ring-LWE and Ideal Lattice Cryptography}\\[2pt]
  {\large\bfseries Ring-LWE 與理想格密碼學}\\[1pt]
  {\small\itshape\color{accent}
    幾何數論、分圓域、Dedekind 整環與後量子安全基礎}\\[5pt]
  \rule{\linewidth}{0.8pt}\\[2pt]
  {\small 黃崇晉（Chung-Chin Huang）\;|\; L16141149
    \;|\; 數論（一）\;|\; 2026 春季，國立成功大學}
\end{center}

\vspace{1pt}

%% ──────────────────────────────────────────────────────────────────────
\section{Section 1 \;—\; 幾何數論基礎：格與 Minkowski 定理
  \textnormal{\small\color{gray}(Geometry of Numbers)}}

\begin{definition}[格（Lattice）]
  取 $\mathbb{R}^n$ 中 $n$ 個線性獨立向量 $b_1,\ldots,b_n$，其整數線性組合的集合
  $\mathcal{L}=\sum_{i=1}^n\mathbb{Z}\,b_i\subset\mathbb{R}^n$
  稱為秩 $n$ 的\textbf{格}。
  基底矩陣 $B=[b_1|\cdots|b_n]$ 的行列式絕對值
  $\det(\mathcal{L}):=|\det B|$ 稱為\textbf{覆積}，與基底選擇無關。
\end{definition}

先證 Blichfeldt 引理作為工具，再推導 Minkowski 定理，最後給出在代數數論中的應用。

\begin{lemma}[Blichfeldt 1914]
  設 $\mathcal{L}\subset\mathbb{R}^n$ 為格，$S\subset\mathbb{R}^n$ 為 Lebesgue 可測集。
  若 $\mathrm{vol}(S)>\det(\mathcal{L})$，則存在相異 $x,y\in S$ 使得 $x-y\in\mathcal{L}$。
\end{lemma}

\begin{proofbox}[證明（引理 1）]
  設 $\mathcal{F}$ 為 $\mathcal{L}$ 的基本域，$\mathrm{vol}(\mathcal{F})=\det(\mathcal{L})$。
  令 $S_v:=S\cap(v+\mathcal{F})$ 並令 $S_v':=S_v-v\subset\mathcal{F}$；
  $\{S_v\}_{v\in\mathcal{L}}$ 為 $S$ 的不相交分解，故
  $\sum_{v\in\mathcal{L}}\mathrm{vol}(S_v')=\mathrm{vol}(S)>\mathrm{vol}(\mathcal{F})$。
  由鴿巢原理，存在相異 $v_1,v_2\in\mathcal{L}$ 使 $S_{v_1}'\cap S_{v_2}'\neq\varnothing$；
  取 $z$ 於此交集，令 $x=z+v_1,\,y=z+v_2\in S$，
  則 $x-y=v_1-v_2\in\mathcal{L}\setminus\{0\}$。
\end{proofbox}

\begin{theorem}[Minkowski 第一定理 1896]
  設 $\mathcal{L}\subset\mathbb{R}^n$ 為格，$S\subset\mathbb{R}^n$ 為中心對稱凸體
  （$x\in S\Rightarrow{-x}\in S$，$S$ 有界閉凸）。
  若 $\mathrm{vol}(S)>2^n\det(\mathcal{L})$，則 $S$ 含至少一個非零格點。
\end{theorem}

\begin{proofbox}[證明（定理 1）]
  令 $\tfrac{1}{2}S=\{\tfrac{1}{2}x:x\in S\}$，則
  $\mathrm{vol}(\tfrac{1}{2}S)=\tfrac{\mathrm{vol}(S)}{2^n}>\det(\mathcal{L})$。
  由引理 1，存在相異 $p,q\in\tfrac{1}{2}S$ 使 $p-q\in\mathcal{L}$。
  由 $2p,2q\in S$，中心對稱知 $-2q\in S$，凸性得
  $p-q=\tfrac{2p+(-2q)}{2}\in S$；故 $p-q\in S\cap\mathcal{L}\setminus\{0\}$。
\end{proofbox}

\textbf{應用：Minkowski 界與理想存在性。}
設 $K$ 為次數 $n=r_1+2r_2$ 的數域，$\Delta_K$ 為判別式，
$\mathcal{O}_K$ 為整數環。透過標準嵌入 $\sigma:K\hookrightarrow\mathbb{R}^{r_1}\times\mathbb{C}^{r_2}\cong\mathbb{R}^n$，
每個理想 $\mathfrak{a}\subset\mathcal{O}_K$ 給出格 $\sigma(\mathfrak{a})$。
對乘積條件 $\prod_i|x_i|^{n_i}=1$（$n_i=1$ 或 $2$）的中心對稱凸體套用定理 1，可得：

\begin{corollary}[Minkowski 界]
  $\mathrm{Cl}(K)$ 中每個理想類均含範數有界的整數理想 $\mathfrak{a}$：
  \[
    N(\mathfrak{a})\;\leq\;
    M_K\;:=\;\sqrt{|\Delta_K|}\cdot
    \Bigl(\tfrac{4}{\pi}\Bigr)^{r_2}\cdot\frac{n!}{n^n}.
  \]
  此式保證 $M_K<\infty$，從而 $\mathrm{Cl}(K)$ 為有限群（類數有限）。
\end{corollary}

Minkowski 定理的\textbf{非構造性}連結幾何（格體積）與算術（理想範數）：
它保證格點\emph{存在}，卻不給出尋找方法——這既是 SVP 困難性的根源，也是密碼學安全性的幾何基石。

%% ──────────────────────────────────────────────────────────────────────
\section{Section 2 \;—\; 分圓域與理想格結構
  \textnormal{\small\color{gray}(Cyclotomic Fields and Ideal Lattices)}}

\textbf{分圓域設定。}
令 $n=2^k$，$\zeta=\zeta_{2n}=e^{i\pi/n}$ 為第 $2n$ 個本原單位根，
$\Phi_{2n}(x)=x^n+1$（第 $2n$ 個分圓多項式，$n=2^k$ 時在 $\mathbb{Q}$ 上不可約），
$K=\mathbb{Q}(\zeta)$，$[K:\mathbb{Q}]=n$，$r_1=0$，$r_2=n/2$。

\begin{definition}[Dedekind 整環]
  $\mathcal{O}_K=\mathbb{Z}[\zeta]$（$K$ 的代數整數環）是 \textbf{Dedekind 整環}：
  Noetherian、整閉，且每個非零質理想皆為極大理想。
  由此，$\mathcal{O}_K$ 中每個非零理想均有唯一的質理想積分解：
  $\mathfrak{a}=\mathfrak{p}_1^{e_1}\cdots\mathfrak{p}_r^{e_r}$（質理想唯一分解性）。
\end{definition}

\begin{definition}[標準嵌入（Canonical Embedding）]
  $K$ 有 $n$ 個複數嵌入 $\sigma_j:K\to\mathbb{C}$，定義為 $\sigma_j(\zeta)=\zeta^{2j-1}$，$j=1,\ldots,n$。
  \textbf{標準嵌入}為
  \[
    \sigma\colon K\;\longrightarrow\;\mathbb{C}^n,\quad
    \alpha\;\longmapsto\;(\sigma_1(\alpha),\ldots,\sigma_n(\alpha)).
  \]
  因共軛對稱 $\sigma_{n+1-j}(\alpha)=\overline{\sigma_j(\alpha)}$，
  $\sigma$ 的像落在同構於 $\mathbb{R}^n$ 的子空間 $H\subset\mathbb{C}^n$ 中。
  標準嵌入使環乘法與 $\mathbb{C}^n$ 中的逐分量乘法相容，
  是連結代數結構與格幾何的橋樑。
\end{definition}

\textbf{判別式與格行列式。}
$\sigma(\mathcal{O}_K)\subset H$ 是秩 $n$ 的格，其覆積由\textbf{判別式} $\Delta_K$ 決定：
\[
  \det\bigl(\sigma(\mathcal{O}_K)\bigr)\;=\;\sqrt{|\Delta_K|}.
\]
對 $K=\mathbb{Q}(\zeta_{2n})$（$n=2^k$），利用分圓域判別式公式
$|\Delta_K|=\tfrac{(2n)^n}{2^n}=n^n$，故 $\det(\sigma(\mathcal{O}_K))=n^{n/2}$。

\begin{definition}[理想格與範數]
  $\mathcal{O}_K$ 的分式理想 $\mathfrak{a}$ 透過標準嵌入給出
  $\sigma(\mathfrak{a})\subset H$，稱為\textbf{理想格}，秩 $n$，且在 $\mathcal{O}_K$ 的乘法作用下封閉。
  理想範數 $N(\mathfrak{a})=[\mathcal{O}_K:\mathfrak{a}]$ 與覆積的關係為
  \[
    \det\bigl(\sigma(\mathfrak{a})\bigr)\;=\;N(\mathfrak{a})\cdot\sqrt{|\Delta_K|},
  \]
  因為 $\sigma(\mathfrak{a})$ 是 $\sigma(\mathcal{O}_K)$ 的指數 $N(\mathfrak{a})$ 子格。
\end{definition}

%% ──────────────────────────────────────────────────────────────────────
\section{Section 3 \;—\; Ring-LWE 的數論結構
  \textnormal{\small\color{gray}(Number-Theoretic Structure of Ring-LWE)}}

\textbf{問題設定。}令 $R=\mathcal{O}_K=\mathbb{Z}[x]/\Phi_{2n}(x)$，$q$ 為正整數，
商環 $R_q=R/qR\cong\mathbb{Z}_q[x]/\Phi_{2n}(x)$。

\begin{definition}[Ring-LWE（LPR 2010）]
  對秘密 $s\in R_q$ 與誤差分布 $\chi_R$（$R$ 上的離散 Gaussian），
  \textbf{Ring-LWE 假設}稱：分布
  $(a,\;a\cdot s+e)$（$a\overset{\$}{\leftarrow}R_q$，$e\leftarrow\chi_R$）
  與 $R_q\times R_q$ 上均勻分布計算不可區分。
  Lyubashevsky--Peikert--Regev（2010）建立量子歸約：
  worst-case ideal-SVP$_\gamma\;\to\;$average-case Ring-LWE。
\end{definition}

\textbf{質理想分解與 $R_q$ 的結構。}
對質數 $q\nmid 2n$，$q$ 在 $\mathcal{O}_K$ 中的分解由 $\Phi_{2n}(x)\bmod q$ 的因式分解決定：
\[
  \Phi_{2n}(x)\;\equiv\;\prod_{i=1}^{g}\phi_i(x)^{e_i}\pmod{q},\quad
  \deg\phi_i=f,\quad e\cdot f\cdot g=n.
\]
對應地 $q\mathcal{O}_K=\mathfrak{q}_1^{e_1}\cdots\mathfrak{q}_g^{e_g}$，
各質理想範數 $N(\mathfrak{q}_i)=q^f$。\textbf{完全分裂}（$e=1,\,f=1,\,g=n$）當且僅當
$q\equiv 1\pmod{2n}$，此時 $\mathbb{F}_q^\times$ 含 $2n$ 次本原單位根 $\omega$，
$\Phi_{2n}(x)\equiv\prod_{j=1}^n(x-\omega_j)\pmod{q}$，由中國剩餘定理：
\[
  R_q\;\cong\;\mathbb{F}_q^n\quad(\text{環同構}).
\]

\begin{definition}[NTT 的數論根源]
  完全分裂條件使 $R_q$ 的乘法化為 $\mathbb{F}_q^n$ 中的逐分量乘法：
  \[
    (a\cdot b)\bmod\Phi_{2n}(x)
    \;\longleftrightarrow\;
    \bigl(\hat{a}(\omega_j)\cdot\hat{b}(\omega_j)\bigr)_{j=1}^n.
  \]
  評估映射 $a\mapsto(\hat{a}(\omega_1),\ldots,\hat{a}(\omega_n))$ 即為\textbf{數論變換（NTT）}，
  利用 $\omega_j$ 的根號結構以 Cooley-Tukey 蝴蝶運算可在 $O(n\log n)$ 時間完成，
  遠優於直接多項式乘法的 $O(n^2)$。
\end{definition}

\textbf{差異理想與逆差理想。}
$\mathcal{O}_K$ 對 $\mathbb{Q}$ 的\textbf{差異理想}（different）由最小多項式的導數生成：
\[
  \mathfrak{D}_{K/\mathbb{Q}}\;=\;\bigl(\Phi_{2n}'(\zeta)\bigr)
  \;=\;\bigl(n\zeta^{n-1}\bigr)
  \;=\;(n),
\]
因 $\zeta^{n-1}=-\zeta^{-1}$ 為 $\mathcal{O}_K$ 中的可逆元（單位根）。
\textbf{逆差理想}（codifferent）為跡型式意義下 $\mathcal{O}_K$ 的對偶格：
\[
  \mathfrak{D}^{-1}
  \;=\;\bigl\{x\in K:\mathrm{Tr}_{K/\mathbb{Q}}(x\,\mathcal{O}_K)\subset\mathbb{Z}\bigr\}
  \;=\;\tfrac{1}{n}\,\mathcal{O}_K.
\]

\begin{proofbox}[$\mathfrak{D}^{-1}=\tfrac{1}{n}\mathcal{O}_K$ 的驗證]
  $\{1,\zeta,\ldots,\zeta^{n-1}\}$ 為 $\mathcal{O}_K$ 的 $\mathbb{Z}$-基。
  利用字元求和，對 $0\le j\le n-1$：
  \[
    \mathrm{Tr}_{K/\mathbb{Q}}(\zeta^j)
    \;=\;\sum_{k\text{ odd}}e^{i\pi jk/n}
    \;=\;e^{i\pi j/n}\sum_{k=0}^{n-1}e^{2\pi ijk/n}
    \;=\;\begin{cases}n&j=0,\\0&1\le j\le n-1.\end{cases}
  \]
  故對 $x=\frac{1}{n}\sum c_j\zeta^j\in\frac{1}{n}\mathcal{O}_K$ 與任意
  $y=\sum a_k\zeta^k\in\mathcal{O}_K$：
  $\mathrm{Tr}(xy)=\frac{1}{n}\mathrm{Tr}((\sum c_j\zeta^j)(\sum a_k\zeta^k))=\frac{n}{n}\cdot(\text{常數項})\in\mathbb{Z}$。
  逆包含由 $\mathrm{Tr}(\frac{1}{m}\cdot 1)=\frac{n}{m}\in\mathbb{Z}$ 要求 $m\mid n$ 得出。
\end{proofbox}

\textbf{密碼學意義。}Ring-LWE 的嚴格表述中，誤差的自然定義域為
$\mathfrak{D}^{-1}/q\mathfrak{D}^{-1}=\frac{1}{n}R_q$（而非 $R_q$）。
這是因為跡配對 $\langle a,b\rangle=\mathrm{Tr}_{K/\mathbb{Q}}(ab)$ 使
$\mathfrak{D}^{-1}$ 恰為 $\mathcal{O}_K$ 在此配對下的\textbf{對偶格}，
確保誤差分布的代數不變性與安全歸約的嚴密性。

%% ──────────────────────────────────────────────────────────────────────
\section{Section 4 \;—\; 應用與延伸
  \textnormal{\small\color{gray}(Applications and Open Problems)}}

\textbf{NIST 後量子標準（2024 年 8 月）。}
\begin{itemize}
  \item \textbf{FIPS 203（ML-KEM $=$ Kyber）}：金鑰封裝，基於 Module-LWE
    （Ring-LWE 在模組格 $R_q^{k\times k}$ 上的推廣）；
    Kyber-512 使用 $n=256$，$q=3329$，NIST Level I（攻擊難度 $\approx 2^{118}$）。
  \item \textbf{FIPS 204（ML-DSA $=$ Dilithium）}：數位簽章，
    基於 Module-LWE 與 Module-SIS；
    安全性最終歸約至理想格上的 SVP/CVP 困難性。
\end{itemize}
兩者的安全鏈為：破解密碼 $\Leftrightarrow$ 解 BDD on ideal lattice $\Leftarrow$ ideal-SVP$_\gamma$ 困難。

\textbf{開放問題。}
\begin{itemize}
  \item \textbf{Ideal-SVP vs SVP}：理想格的 SVP 是否真的不比任意格容易？
    Cramer et al.（2016）對特定分圓環給出量子加速，但一般情形仍開放；
    若 Ideal-SVP 易解，則 Ring-LWE 安全保證將弱化。
  \item \textbf{非分圓域的 Ring-LWE}：
    LPR（2010）的歸約依賴分圓域的交換 Galois 結構；
    對 Galois 閉包為非交換群的數環，安全歸約是否成立仍未知。
  \item \textbf{古典 worst-case 歸約}：
    Regev 的 LWE 歸約為量子的；
    純古典 worst-case $\to$ average-case 歸約至今未找到。
  \item \textbf{SVP 量子下界}：
    最佳量子演算法為 $2^{0.265n}$（Laarhoven 2015），
    但無條件指數下界至今未被證明。
\end{itemize}

\vspace{4pt}
\noindent\rule{\linewidth}{0.4pt}
\begin{center}\small\color{gray}
  O.\ Regev, \textit{J.\ ACM} \textbf{56}(6), 2009.\;
  V.\ Lyubashevsky--C.\ Peikert--O.\ Regev, \textit{J.\ ACM} \textbf{60}(6), 2013.\\[1pt]
  H.\ Minkowski, \textit{Geometrie der Zahlen}, 1896.\;
  H.\ Blichfeldt, \textit{Trans.\ AMS} \textbf{15}(3), 1914.\\[1pt]
  J.\ Neukirch, \textit{Algebraic Number Theory}, Springer, 1999.\;
  NIST FIPS 203/204, 2024.
\end{center}

\end{document}
